% ==============================================================================
% SEÇÃO 4: USER STORIES DE ACORDO COM CADA FASE E RASTREABILIDADE
% Autores: Alêkson (#4), Italo (#6), Matheus (#8)
%
% REGRA DE OURO ANTI-REDUNDÂNCIA:
% Não reexplicar a matemática de ML nesta seção. O foco aqui é estritamente 
% em personas, declarações ágeis de valor e cenários de teste BDD (Dado-Quando-Então).
% ==============================================================================

\section{User Stories de Acordo com Cada Fase}
\label{sec:user_stories}

Este capítulo formaliza a visão de produto e os requisitos de negócio sob a perspectiva ágil, estruturando as Histórias de Usuário vinculadas a cada fase evolutiva do sistema e estabelecendo os critérios de aceitação em formato BDD (\textit{Behavior-Driven Development}).

% ------------------------------------------------------------------------------
% 4.1 FASE 1: VISÃO DE MERCADO
% Autor: Alêkson Castro (Issue #4)
% ------------------------------------------------------------------------------
\subsection{User Story 1: Visão de Mercado (Fase 1)}
\label{sec:us1_mercado}

% [PLACEHOLDER - ALÊKSON: Redigir a US1 em padrão BDD sintético]
% - Candidato explorando mercado: declaração oficial e cenários BDD de modalidades e salários;
% - Referenciar a Seção 2 (\ref{fig:fluxo_aba1}) para detalhes de tela.
\noindent\textit{[Espaço reservado para a User Story 1 por Alêkson Castro -- Issue \#4. Foco na declaração de valor ágil e nos cenários BDD da visão de mercado.]}

% ------------------------------------------------------------------------------
% 4.2 FASE 1: CBF E COLD START
% Autor: Alêkson Castro (Issue #4)
% ------------------------------------------------------------------------------
\subsection{User Story 2: Filtragem por Conteúdo e Cold Start (Fase 1)}
\label{sec:us2_cold_start}

% [PLACEHOLDER - ALÊKSON: Redigir a US2 em padrão BDD sintético]
% - Candidato sem histórico: declaração oficial e cenários BDD de busca por skills e scoped rerun;
% - Referenciar a Seção 2 (\ref{fig:fluxo_aba2}) para detalhes de tela.
\noindent\textit{[Espaço reservado para a User Story 2 por Alêkson Castro -- Issue \#4. Foco no cold start e nos cenários BDD da busca por competências.]}

% ------------------------------------------------------------------------------
% 4.3 FASE 2: FILTRAGEM COLABORATIVA E EXPLICABILIDADE
% Autor: Italo Emannuel Beckman (Issue #6)
% ------------------------------------------------------------------------------
\subsection{Fase 2: Filtragem Colaborativa e Explicabilidade (US3 e US4)}
\label{sec:us_fase2}

% [PLACEHOLDER - ITALO: Redigir US3 e US4 em padrão BDD sintético]
% - US3 (Recrutador / Perfil latente): Declaração oficial + Cenários BDD de persona e catálogo sem repetição;
% - US4 (Explicabilidade / Waterfall): Declaração oficial + Cenários BDD do gráfico Waterfall Altair;
% - Referenciar a Seção 2 (\ref{fig:fluxo_aba3}) para detalhes de tela.
\noindent\textit{[Espaço reservado para as User Stories 3 e 4 por Italo Beckman -- Issue \#6. Foco na transparência da IA e cenários BDD da CF.]}

% ==============================================================================
% 4.4 USER STORY - FASE 3 (O AVALIADOR DO PROJETO)
% Autor: Matheus Alexandre Nardi (Issue #8)
% Finalidade: Formalização da US5 no padrão BDD e consolidação da Matriz Geral
% ==============================================================================

\subsection{Fase 3: O Avaliador do Projeto (US5)}
\label{sec:us5_avaliador}

Enquanto as histórias anteriores abordam a experiência operacional do candidato e da plataforma de recrutamento, a quinta História de Usuário atende à perspectiva do \textbf{Avaliador do Projeto} (representado pela Professora da disciplina e pela banca examinadora). Sob a ótica do TDSP (\textit{Team Data Science Process}), o papel do avaliador técnico consiste em auditar o rigor estatístico das métricas reportadas, inspecionar a conformidade metodológica e julgar se a complexidade de engenharia da arquitetura híbrida se justifica frente a modelos puros e pisos teóricos de ruído.

\begin{tcolorbox}[colback=blue!5!white, colframe=blue!75!black, title={\textbf{User Story 5 (Obrigatória pelo Roteiro da Disciplina)}}]
\textit{``Como avaliador do projeto, quero comparar as três técnicas lado a lado com métricas reais e inspecionar a metodologia de testes, para julgar se a complexidade da hibridização se justificou perante baselines e pisos teóricos.''}
\vspace{0.15cm}
\tcblower
\textbf{Persona:} Professora da Disciplina / Avaliador Técnico de Sistemas de Recomendação.\\
\textbf{Dor Central de Negócio:} Vulnerabilidade a modelos que funcionam como ``caixas-pretas'', relatórios sem significância estatística comprovada ou soluções superengenheiradas cujos ganhos preditivos não compensam o custo operacional em produção.
\end{tcolorbox}

\subsubsection{Critérios de Aceite Formais (Padrão BDD)}

A homologação da US5 é delimitada por 3 cenários operacionais no padrão BDD (\textit{Behavior-Driven Development}):

\paragraph{Cenário 1: Auditoria do Duelo de Modelos e Piso Irredutível de Ruído.}
\begin{itemize}[leftmargin=1.5em, itemsep=2pt]
    \item \textbf{Dado que} o avaliador acessa o \textit{Modo Avaliador $\to$ Duelo dos Modelos} no painel interativo da aplicação;
    \item \textbf{Quando} inspeciona o quadro comparativo consolidado pelas 8 abordagens técnicas (chutes triviais, baseline de popularidade, oráculo, CBF quantitativa, KNN e SVD);
    \item \textbf{Então} o sistema deve comprovar empiricamente a superioridade do modelo SVD ($\text{RMSE} = 0{,}625$ e $\text{Precision@10} = 1{,}000$) contra o KNN ($\text{RMSE} = 1{,}263$ e $\text{Precision@10} = 0{,}874$) com significância atestada pelo teste de Wilcoxon ($p < 0{,}05$), evidenciando que o erro do SVD situa-se a apenas $1{,}20\times$ do piso teórico bayesiano do oráculo ($\text{RMSE}_{\text{oráculo}} = 0{,}5214$).
    \item \textit{Rastreabilidade Arquitetural:} Fluxo 4 da Arquitetura Alvo (\ref{fig:fluxo_aba4}).
\end{itemize}

\paragraph{Cenário 2: Transparência Metodológica, LGPD e Proveniência de Dados.}
\begin{itemize}[leftmargin=1.5em, itemsep=2pt]
    \item \textbf{Dado que} o avaliador acessa a página \textit{Como avaliamos \& limitações} (Fluxo 5);
    \item \textbf{Quando} consulta o painel de governança e integridade de artefatos;
    \item \textbf{Então} o sistema deve discriminar a data/hora e o módulo gerador de cada artefato serializado de métricas, comprovar a conformidade com a LGPD pelo uso de dados corporativos públicos e personas sintéticas sem qualquer PII, e declarar explicitamente a natureza \textit{self-fulfilling} da avaliação bilinear e as limitações de cobertura textual do catálogo.
\end{itemize}

\paragraph{Cenário 3: Julgamento Técnico da Necessidade de Hibridização.}
\begin{itemize}[leftmargin=1.5em, itemsep=2pt]
    \item \textbf{Dado que} o avaliador examina os trade-offs inerentes às abordagens isoladas (a hiperespecialização semântica da CBF versus o colapso por \textit{cold start} na CF);
    \item \textbf{Quando} confronta a arquitetura alvo da Fase 3 com as alternativas puras;
    \item \textbf{Então} o sistema deve evidenciar que a conjugação de chaveamento dinâmico (\textit{Switching} para $N_{\text{interações}} < 5$) com fusão ponderada ($0{,}4\,\text{CBF} + 0{,}6\,\text{CF}$) sobre escore contínuo bruto é a única estratégia capaz de assegurar cobertura imediata a novos usuários e serendipidade comportamental a candidatos ativos.
    \item \textit{Rastreabilidade Arquitetural:} Fluxo Híbrido da Fase 3 (\ref{fig:fluxo_hibrido_fase3}).
\end{itemize}

% ==============================================================================
% 4.5 MATRIZ GERAL DE RASTREABILIDADE
% ==============================================================================
\subsection{Matriz Geral de Rastreabilidade do Sistema}
\label{sec:matriz_rastreabilidade}

Para assegurar o fechamento metodológico do TDSP e conectar formalmente o valor de negócio à engenharia de software desenvolvida, a Tabela~\ref{tab:matriz_rastreabilidade} sintetiza o encadeamento integral entre as 5 Histórias de Usuário, as personas centrais, os módulos funcionais do sistema, os modos de navegação do dashboard e as métricas de validação empírica.

\begin{table}[H]
\centering
\small
\renewcommand{\arraystretch}{1.25}
\caption{Matriz Geral de Rastreabilidade: Conexão entre User Stories, Personas, Módulos e Validação.}
\label{tab:matriz_rastreabilidade}
\resizebox{\textwidth}{!}{%
\begin{tabular}{@{}llllll@{}}
\toprule
\textbf{User Story} & \textbf{Persona Central} & \textbf{Fase} & \textbf{Modo no Dashboard} & \textbf{Componente / Módulo} & \textbf{Métrica de Validação} \\ \midrule
\textbf{US 1 (Mercado)} & Candidato em Recolocação & 1 & Modo Candidato $\to$ EDA (Tabs H1--H5) & Módulo de EDA \& Limpeza & Medianas Salariais e Salto de CTR \\
\textbf{US 2 (Cold Start)} & Candidato sem Histórico & 1 & Modo Candidato $\to$ Simulador CBF & Motor CBF (\textit{RecSysCBF}) & Cosseno $\ge 0{,}50$ e Re-ranking Scoped \\
\textbf{US 3 (Perfil Latente)} & Plataforma / Recrutador & 2 & Modo Candidato $\to$ Feed CF (SVD) & Motor CF (SVD Latente) & Precision@10 ($1{,}000$) e Escore Bruto \\
\textbf{US 4 (Explicabilidade)} & Candidato com Histórico & 2 & Modo Candidato $\to$ Waterfall Altair & Explicabilidade CF (Altair) & Decomposição Aditiva $\mu + b_u + b_i + q_i^T p_u$ \\
\textbf{US 5 (Avaliador)} & Professora / Avaliador & 3 & Modo Avaliador $\to$ Duelo e Limitações & Pipeline de Benchmarking & Wilcoxon ($p < 0{,}05$), Oráculo e P@10 CBF \\ \bottomrule
\end{tabular}%
}
\end{table}

A matriz evidencia que cada tela da aplicação atende a uma dor concreta de negócio mapeada desde a concepção do produto, fornecendo um roteiro verificável para a homologação dos critérios de aceite detalhados na Seção~6.
